% =====================================================================
%  Annexes
% =====================================================================
%  Le guide demande que les details techniques soient reportes en annexe
%  et que le code source ne figure pas dans le corps. Les annexes sont
%  volontairement breves : le depot du projet porte l'integralite des
%  elements, et sa reference figure ci-dessous.
% =====================================================================
\cleardoublepage
\phantomsection
\addcontentsline{toc}{chapter}{Annexes}
\chapter*{Annexes}
\markboth{Annexes}{}

% ---------------------------------------------------------------------
\section*{Annexe A -- Depot du projet et mise en route}
\addcontentsline{toc}{section}{Annexe A -- Depot du projet et mise en route}

L'integralite du code source, des modeles UML et de la documentation
technique est publiee a l'adresse suivante :

\begin{center}
\url{\depotGitHub}
\end{center}

Le depot contient notamment deux cent vingt modeles UML au format
PlantUML, repartis en vues de plateforme, d'infrastructure, de frontal et
de microservice. Ces modeles sont du texte : ils se versionnent et se
comparent comme du code, et c'est de leur rendu que proviennent toutes
les figures de conception de ce rapport.

\paragraph{Mise en route.}
La plateforme se lance en une commande, depuis la racine du depot :
\texttt{docker compose up -{}-build}. Les vingt conteneurs demarrent dans
l'ordre impose par leurs dependances, la base et le courtier de messages
d'abord. Le back-office est alors accessible sur le port 3000, le portail
client sur le port 3001, la passerelle sur le port 8080.

% ---------------------------------------------------------------------
\section*{Annexe B -- Extrait du journal de bord}
\addcontentsline{toc}{section}{Annexe B -- Extrait du journal de bord}

\begin{table}[H]
\centering
\small
\begin{tabularx}{\textwidth}{@{}p{2.3cm}X@{}}
\toprule
\textbf{Periode} & \textbf{Faits marquants et enseignements} \\ \midrule

Sprint 1 & La suppression d'un compte est un faux besoin : les comptes
  portent un historique metier. L'archivage reversible a ete introduit
  apres avoir constate qu'une suppression cassait les references \\

Sprint 2 & La surveillance des echeances a impose de stocker les dates
  plutot que de les deduire ; la portee de ce choix n'est apparue qu'au
  sprint~5 \\

Sprint 3 & Le point d'entree creant une mission sans devis a ete retire :
  il contournait la tarification. Difficulte : separer la negociation de
  l'execution \\

Sprint 4 & Le huitieme service a ete ajoute sans modifier aucun service
  existant. Choix assume d'un moteur de regles explicables plutot que
  d'un modele appris sans donnees \\

Sprint 5 & Divergence entre l'enumeration du service d'alertes et celle
  du projet commun, qui faisait echouer la lecture des donnees \\

Sprint 6 & La liaison permanente de la messagerie a impose d'adapter la
  passerelle, seul endroit portant une regle de proxy specifique \\

Recette & Mise aux normes des modeles UML, jeu de donnees de
  demonstration, redaction du rapport \\

\bottomrule
\end{tabularx}
\end{table}

% ---------------------------------------------------------------------
\section*{Annexe C -- Jeu de donnees de demonstration}
\addcontentsline{toc}{section}{Annexe C -- Jeu de donnees de demonstration}

Les captures d'ecran de ce rapport ont ete prises sur une base peuplee
par un script idempotent, verse au depot. Ce jeu de donnees couvre
volontairement \essentiel{tous les etats de chaque cycle de vie}, afin que
les ecrans presentes soient representatifs et non anecdotiques.

\begin{table}[H]
\centering
\small
\begin{tabularx}{\textwidth}{@{}p{4.3cm}X@{}}
\toprule
\textbf{Module} & \textbf{Volumetrie et couverture} \\ \midrule
Utilisateurs & 42 comptes, dont 4 en attente et 15 archives ; six roles \\
Flotte & 16 vehicules sur six statuts, 14 equipements, 18 pieces dont 3
  sous leur seuil \\
Missions & 17 demandes couvrant les cinq statuts, 32 missions, 69
  positions sur quatre trajets \\
Maintenance & 17 ordres de travail sur les quatre colonnes, 18
  interventions \\
Finance et alertes & 17 factures sur les cinq statuts, 46 alertes des
  quatre priorites \\
Messagerie & 12 conversations, 47 messages dont deux pieces jointes \\
\bottomrule
\end{tabularx}
\end{table}

Les villes, distances et coordonnees employees sont reelles, les
immatriculations suivent le format tunisien et les montants sont en
dinars. Les avatars sont des monogrammes generes : aucune photographie de
personne n'a ete utilisee.
